% !TeX root = ../main.tex

\section{Detection Methods}
\label{sec:methods}

The detector examines an incoming training dataset before the continual
learner is updated. It asks whether the images contain evidence of a
forgetting attack; it does not use the detection decision to change the
learner in the experiments reported here. Two information settings are
considered. The supervised detector learns from clean and BrainWash-poisoned
images on an earlier task. The reference-based methods instead compare the
incoming data with synthetic samples obtained by inverting historical
models. They do not fit a clean/poison classifier.

Both settings use responses of a frozen model. They differ in what is learned
from those responses and in how a dataset decision is made. In particular,
Rank is not the aggregation rule for the supervised detector: the supervised
dataset decision is obtained by counting image-level logistic-regression
alerts. Rank performs a separate, direct comparison of dataset-level feature
correlations.

\subsection{Features of a Frozen Continual Learner}
\label{subsec:methods-features}

\subsubsection{The incoming head as a fixed probe}

Tasks are indexed from zero. Before incoming task $t$, the backbone
$g_{\theta_{<t}}$ has learned the preceding clean tasks
$0,\ldots,t-1$. An untrained head $h_t$ is attached to this backbone for the
$K$ incoming classes. Its parameters are initialised with a fixed seed and
kept unchanged during extraction. The predictive vector is
\begin{equation}
    p_t(x)=\operatorname{softmax}
    \bigl(h_t(g_{\theta_{<t}}(x))\bigr),
    \qquad
    \widehat{k}_t(x)=\operatorname*{argmax}_{1\leq k\leq K}p_{t,k}(x).
    \label{eq:methods-probe}
\end{equation}
The head has not learned to classify the incoming images. Its output is a
repeatable probe of the current representation, rather than an estimate of
the image's correct class. Reusing its initialisation prevents differences
between clean and modified views from being caused by different random
heads. It does not eliminate differences between backbones trained on
different histories.

For each image, the extraction loss is
\begin{equation}
    \ell_t(x)=-\log p_{t,\widehat{k}_t(x)}(x).
    \label{eq:methods-predicted-loss}
\end{equation}
The predicted class is treated as a fixed target for this backward pass.
The true image-class label is not read. This is different from taking the
gradient of cross-entropy with respect to the true class and merely removing
that label from the saved feature table. The loss in
Equation~\eqref{eq:methods-predicted-loss} also contains no continual-learning
regularisation term.

The model remains in evaluation mode. Gradients are calculated without an
optimiser step, and neither the parameters nor the BatchNorm running
statistics are updated. Each image is processed separately, so its features
do not depend on the other members of a simulated dataset.

\subsubsection{Parameter-tensor gradient norms}

Let $\Theta_1,\ldots,\Theta_L$ denote the backbone's parameter tensors,
excluding every task head. The gradient feature associated with tensor
$\Theta_\ell$ is
\begin{equation}
    v_{t,\ell}(x)=
    \left\|\nabla_{\Theta_\ell}\ell_t(x)\right\|_{\mathrm{F}},
    \qquad \ell=1,\ldots,L.
    \label{eq:methods-tensor-norm}
\end{equation}
The Frobenius norm here is the Euclidean norm after flattening the tensor.
One convolutional weight tensor therefore supplies one feature, not one
feature per scalar weight. Biases and the trainable BatchNorm scale and
offset tensors are retained separately. In the implemented backbone this
gives $L=80$ features.

These norms preserve where gradient magnitude occurs in the network while
discarding directions within a tensor. A coarser description combines
tensors into the stem and four residual stages. If $\mathcal{S}_q$ is the
set of tensors in stage $q$, its norm is
\begin{equation}
    v^{\mathrm{stage}}_{t,q}(x)=
    \left(\sum_{\ell\in\mathcal{S}_q}v_{t,\ell}(x)^2\right)^{1/2}.
    \label{eq:methods-stage-norm}
\end{equation}
The corresponding global norm combines all five stages in the same way.
This relation makes the stage-level ablation a controlled reduction in
granularity, rather than a change of gradient target or model state.

\subsubsection{Prediction and activation features}

Four additional features describe the same frozen forward pass. Writing
$p_{t,(1)}(x)$ and $p_{t,(2)}(x)$ for the largest and second-largest entries
of $p_t(x)$, these are
\begin{align}
    H_t(x)&=-\sum_{k=1}^{K}p_{t,k}(x)\log p_{t,k}(x),
        &c_t(x)&=p_{t,(1)}(x), \label{eq:methods-uncertainty-a}\\
    m_t(x)&=p_{t,(1)}(x)-p_{t,(2)}(x),
        &a_t(x)&=\|g_{\theta_{<t}}(x)\|_2.
        \label{eq:methods-uncertainty-b}
\end{align}
The activation is the vector entering the incoming head, after pooling and
flattening. Entropy, confidence and margin describe the probe response;
none is assumed to be calibrated for incoming-class accuracy.

The full supervised feature vector is
\begin{equation}
    \varphi_t(x)=
    \bigl(v_{t,1}(x),\ldots,v_{t,80}(x),H_t(x),c_t(x),m_t(x),a_t(x)\bigr)
    \in\mathbb{R}^{84}.
    \label{eq:methods-full-feature}
\end{equation}
True-class probability, the loss itself, and cosine similarities with
historical gradients are not classifier inputs. Some of these quantities
are retained as extraction diagnostics in the implementation. Their
availability in an output file does not mean they belong to the fitted
detector. The supervised score has no inversion-derived input, although the
experimental pipeline still uses inversions to generate attacks and
reference-based diagnostics.

\subsection{Supervised Image Scoring}
\label{subsec:methods-supervised}

The source task supplies pairs $(x_i,x_i^{\mathrm{BW}})$, where
$x_i^{\mathrm{BW}}$ is the BrainWash-modified version of original $x_i$.
The detector label is $z_i=0$ for the clean view and $z_i=1$ for the poisoned
view. This binary label is unrelated to the image-class label used by the
continual learner. Both views of an original stay in the same data split.

For each feature coordinate, a mean $\mu_j$ and standard deviation $s_j$
are fitted using source-training rows only. The transformed input is
\begin{equation}
    \widetilde{\varphi}_{t,j}(x)
       =\frac{\varphi_{t,j}(x)-\mu_j}{s_j},
    \label{eq:methods-scaling}
\end{equation}
with constant coordinates handled by the scaler without division by zero.
The scaler is reused on later tasks. No target mean or variance is
estimated to align the target features with the source.

A logistic regression maps the standardised features to
\begin{equation}
    s_t(x)=\sigma\bigl(b+w^\top\widetilde{\varphi}_t(x)\bigr),
    \qquad \sigma(u)=(1+\exp(-u))^{-1}.
    \label{eq:methods-logistic-score}
\end{equation}
The fit minimises binary cross-entropy with an $\ell_2$ penalty on $w$.
In the implementation's summed-loss convention, the objective is equivalent
to
\begin{equation}
    \frac{1}{2}\|w\|_2^2+
    C\sum_{i\in\mathcal{I}_{\mathrm{fit}}}
    \left[-z_i\log s(x_i)-(1-z_i)\log(1-s(x_i))\right],
    \label{eq:methods-logistic-objective}
\end{equation}
where $C=1$ and the intercept is not penalised. Clean and poisoned views have
equal training weight. Randomly perturbed images are excluded from the fit.

The regression is fitted once for each continual-learning implementation,
seed and feature variant. Its coefficients are not refitted for different
target poisoning fractions or attack settings. A large score indicates
similarity to the source poisoning examples in this representation; it is
not a calibrated probability that an arbitrary deployment image is
poisoned.

\subsection{From Image Scores to a Dataset Alert}
\label{subsec:methods-thresholds}

Two thresholds are needed. A sample threshold $\tau$ marks suspicious
images, and a count threshold $c_\star$ determines whether there are enough
such images to alert on the dataset. Averaging image scores or fitting a
second logistic regression on datasets is not the decision procedure used
in the final experiments.

\subsubsection{Historical sample-threshold calibration}

Let $\mathcal{H}$ be the declared source and development tasks, and let
$V_h$ contain the clean validation images of task $h\in\mathcal{H}$.
For each task, $\tau_h^{Q}$ is chosen so that the observed fraction of
scores at or above it is at most $\alpha_s$. The implementation sorts the
scores and places the threshold just above the appropriate order statistic,
which also handles tied scores conservatively. Here $\alpha_s=0.05$.

The robust historical rule supplements this clean quantile with a cutoff
in logit space. For validation scores clipped to
$[10^{-12},1-10^{-12}]$, define
\begin{align}
    u_h(x)&=\log\frac{\operatorname{clip}(s_h(x))}
                              {1-\operatorname{clip}(s_h(x))},\\
    M_h&=\operatorname{median}_{x\in V_h}u_h(x),
    \qquad
    D_h=1.4826\operatorname{median}_{x\in V_h}|u_h(x)-M_h|.
    \label{eq:methods-mad}
\end{align}
The chosen sample threshold is
\begin{equation}
    \tau=\max_{h\in\mathcal{H}}
          \left\{\tau_h^{Q},\;\sigma(M_h+3D_h)\right\}.
    \label{eq:methods-historical-cutoff}
\end{equation}
Taking the larger within-task cutoff never lowers the empirical clean
quantile. Taking the maximum across tasks favours conservative alerts on
the observed histories. Neither operation establishes that the same clean
score distribution will hold on the next task.

Two simpler rules are retained for comparison: the source clean quantile
alone, and the maximum clean quantile over the historical calibration
tasks. The principal rule in Equation~\eqref{eq:methods-historical-cutoff},
called \texttt{history\_mad} in the implementation, was fixed before the
expanded runs. The development-task test results diagnose that choice;
they do not trigger a switch to whichever rule performs best on the target.

\subsubsection{Count calibration on separate clean originals}

The sample threshold is applied to clean reserve sets that are disjoint from
the validation images used above. Let $J_h$ of the $M_h^{\mathrm{res}}$
reserve images on historical task $h$ exceed $\tau$. A one-sided
Clopper--Pearson bound~\cite{clopper1934binomial} gives
\begin{equation}
    p_h^+=
    \begin{cases}
      1, & J_h=M_h^{\mathrm{res}},\\
      F_{\mathrm{Beta}(J_h+1,M_h^{\mathrm{res}}-J_h)}^{-1}(1-\eta),
          & J_h<M_h^{\mathrm{res}},
    \end{cases}
    \label{eq:methods-cp}
\end{equation}
where $F^{-1}$ denotes the beta quantile. The maximum historical bound,
$p^+=\max_h p_h^+$, is used for count calibration. For dataset size $N$,
\begin{equation}
    c_\star=\min\left\{c\in\{0,\ldots,N+1\}:
         \Pr\{\operatorname{Bin}(N,p^+)\geq c\}\leq\gamma\right\}.
    \label{eq:methods-critical-count}
\end{equation}
The nominal budget is split as $\eta=\gamma=0.025$. The value $N+1$
represents an unattainable cutoff if no count up to $N$ satisfies the tail
criterion.

For an incoming dataset $B=\{x_1,\ldots,x_N\}$, the supervised decision is
\begin{equation}
    C(B)=\sum_{i=1}^{N}\mathbf{1}\{s_t(x_i)\geq\tau\},
    \qquad
    A(B)=\mathbf{1}\{C(B)\geq c_\star\}.
    \label{eq:methods-dataset-decision}
\end{equation}
The count is calculated directly from the frozen image scores. Its
calibration assumes an appropriate binomial model for clean alerts and
stability of the clean alert probability. Class-dependent images, repeated
sampling from a finite pool, and task shift all require caution. In this
study the calculation is a historical calibration heuristic, not a
guarantee of a 5\% target false-positive rate. The empirical target clean
alert rate must be measured separately.

\subsection{Historical-Reference Tests Without Poisoning Labels}
\label{subsec:methods-unsupervised}

The reference-based setting removes access to labelled clean/poison
examples for detector fitting. It assumes instead that the historical
learner and its inversion samples are trustworthy. Synthetic inversion
targets remain available when reconstructing historical images. This
assumption is distinct from knowing which incoming images are clean.

The incoming data and the references are described with the same frozen
backbone and incoming head. These methods use a fixed ten-dimensional
descriptor: the global gradient norm, five stage norms, entropy,
confidence, margin and activation norm. They do not use the supervised
regression's 84-dimensional feature vector, its coefficients or its
thresholds. Thus their comparison with the supervised detector concerns
different information settings and representations, rather than an
aggregation-only ablation.

\subsubsection{MMD on historical descriptors}
\label{subsubsec:methods-mmd}

The baseline measures a discrepancy between an incoming descriptor sample
and a pooled historical reference bank. Maximum mean discrepancy compares
distributions through kernel mean embeddings~\cite{gretton2012mmd}.
Here the kernel is approximated with random Fourier
features~\cite{rahimi2007random}, keeping repeated dataset tests inexpensive.

Historical rows are split into a preprocessing-fit subset and a disjoint
test bank. The fit subset alone determines the transformation: gradient and
activation norms receive a $\log(1+x)$ transform, coordinates are centred
by their medians and scaled by $1.4826$ times their median absolute
deviations, and near-zero-scale coordinates are removed. Standardised
values are clipped to $[-12,12]$. The RBF bandwidth is the median nonzero
pairwise distance in this transformed fit subset. No incoming values are
used to select these settings.

For transformed descriptor $d$, the random map is
\begin{equation}
    \psi(d)=\sqrt{\frac{2}{D}}
       \bigl(\cos(\omega_1^\top d+b_1),\ldots,
             \cos(\omega_D^\top d+b_D)\bigr),
    \label{eq:methods-rff}
\end{equation}
where $\omega_j\sim\mathcal{N}(0,\ell^{-2}I)$ and
$b_j\sim\operatorname{Unif}(0,2\pi)$ are fixed after fitting. The statistic
is the biased squared distance between the two feature means,
\begin{equation}
    T_{\mathrm{MMD}}=
    \left\|\frac{1}{n_R}\sum_{i=1}^{n_R}\psi(d_i^R)
          -\frac{1}{n_B}\sum_{i=1}^{n_B}\psi(d_i^B)\right\|_2^2.
    \label{eq:methods-mmd}
\end{equation}
Group assignments are permuted while preserving group sizes and the frozen
feature map. With $P$ random permutations, the Monte Carlo value is
$(1+\#\{T^{(b)}\geq T_{\mathrm{MMD}}\})/(P+1)$; numerical ties are
counted conservatively. The null concerns equality of the represented
distributions. Its permutation interpretation requires exchangeability,
which is not established merely because the historical checkpoint is
clean. Synthetic inversion data and genuine incoming images can differ
even in the absence of an attack.

\subsubsection{Rank correlations within each dataset}
\label{subsubsec:methods-rank}

Rank compares historical tasks separately rather than pooling their
descriptors. Its aim is to discard changes in individual feature scales
while retaining changes in pairwise dependence. For feature coordinates
$a<b$ and a sample of $n$ rows, the implemented Kendall statistic is
\begin{equation}
    \widehat{\tau}_{ab}
    =\frac{2}{n(n-1)}\sum_{i<j}
       \operatorname{sign}(d_{ia}-d_{ja})
       \operatorname{sign}(d_{ib}-d_{jb}).
    \label{eq:methods-kendall}
\end{equation}
This is Kendall's $\tau_a$~\cite{kendall1938rank}; tied values contribute
zero. It is not a tie-corrected $\tau_b$. Ten descriptors give 45 possible
pairs.

For each historical reference and feature pair, the correlation is also
calculated after leaving out each row. Its jackknife pseudovalue is
\begin{equation}
    P_{i,ab}=n\widehat{\tau}_{ab}-(n-1)\widehat{\tau}_{ab}^{(-i)},
    \qquad
    Q_{i,ab}=P_{i,ab}-\overline{P}_{ab}.
    \label{eq:methods-jackknife}
\end{equation}
These pseudovalues describe first-order sampling variation, not additional
images~\cite{efron1981jackknife}. A pair is retained only if the empirical
variance of its centred pseudovalues exceeds $10^{-12}$ in every historical
reference. The retained set $\mathcal{A}$ is selected using historical
data alone. This excludes effectively degenerate reference coordinates; it
does not prove that all bootstrap assumptions hold.

For incoming dataset $B$ and reference task $h$, the observed discrepancy is
\begin{equation}
    T_h=\max_{(a,b)\in\mathcal{A}}
         |\widehat{\tau}_{ab}^{B}-\widehat{\tau}_{ab}^{h}|.
    \label{eq:methods-rank-statistic}
\end{equation}
A Gaussian multiplier construction approximates its null variation. For
each draw, independent standard-normal row weights are applied to centred
pseudovalues. For a sample of size $n$, the implemented scaling is
\begin{equation}
    G_{ab}^{(b)}=
       \frac{1}{\sqrt{n(n-1)}}\sum_{i=1}^{n}\xi_i^{(b)}Q_{i,ab},
       \qquad \xi_i^{(b)}\sim\mathcal{N}(0,1).
    \label{eq:methods-multiplier}
\end{equation}
The same row weight is shared across feature pairs within a sample, so
their dependence is retained. Incoming and reference weights are generated
separately. The simulated statistic and approximate comparison value are
\begin{align}
    T_h^{(b)}&=\max_{(a,c)\in\mathcal{A}}
        |G_{ac}^{B,(b)}-G_{ac}^{h,(b)}|,\\
    \widehat{p}_h&=
        \frac{1+\sum_{b=1}^{B_0}\mathbf{1}\{T_h^{(b)}\geq T_h\}}
             {B_0+1}.
    \label{eq:methods-rank-pvalue}
\end{align}
There are $B_0=499$ draws, giving a Monte Carlo resolution of $1/500$.
Numerically tied statistics are counted conservatively.
This construction is motivated by Gaussian and jackknife-bootstrap
approximations for nondegenerate high-dimensional
$U$-statistics~\cite{chen2018bootstrap}. Its use here is an experimental
two-sample procedure, not an exact permutation test or a proved
finite-sample poisoning detector.

Rank accepts compatibility with any historical reference. Equivalently,
its null is the union of the taskwise compatibility hypotheses. The dataset
alerts only when every taskwise comparison rejects:
\begin{equation}
    \widehat{p}_{\mathrm{union}}=\max_h\widehat{p}_h,
    \qquad
    A_{\mathrm{Rank}}(B)=
       \mathbf{1}\{\widehat{p}_{\mathrm{union}}\leq 0.05\}.
    \label{eq:methods-rank-union}
\end{equation}
This is not a vote over image classifications. Increasing the number of
historical references gives the incoming data more opportunities to match
a reference, which can reduce sensitivity.

Strictly increasing transformations of individual coordinates leave these
rank correlations unchanged. Rank can therefore miss an attack that alters
only marginal magnitudes. It can also miss higher-order dependence changes
that preserve pairwise correlations. If a retained descriptor becomes
constant in an incoming bag, the implementation reports an unsupported
comparison rather than treating it as evidence that the bag is clean.
Reference fidelity, ties and dependence between inversion samples further
limit the interpretation of the approximate values.

\subsection{What the Decisions Establish}
\label{subsec:methods-interpretation}

For the supervised method, the fitted representation, scaler, regression,
sample cutoff and count cutoff are all frozen before target evaluation.
Target clean/poison labels are used only to construct the controlled
mixtures and calculate metrics. Reference-based fitting likewise precedes
reading the incoming benchmark, and its predictors receive descriptor
columns rather than attack labels or class labels.

An alert has different meanings in the two settings: source-learned
poisoning evidence for logistic regression, and historical distribution
incompatibility for MMD or Rank. None certifies the cause of a discrepancy.
Conversely, a non-alert does not establish that learning the dataset is
safe. Determining whether filtering by these decisions prevents forgetting
would require a further intervention experiment, which is outside the
evaluation described in the next section.
